IMPLEMENT ME; USE A DOUBLE in the callback!!! ///

Feedforward
\begin{DoxyEnumerate}
\item Constant Desired Force\+: Dropping the Feedforward Term\+: Yes, you can potentially drop the feedforward term if you prioritize simplicity and have a robust PID controller. However, be aware of potential tracking errors, overshoot, and longer settling times, especially during transient phases or significant disturbances.
\item Increasingly Changing Setpoint Force\+: Feedforward Term as Next Setpoint\+: This approach can be helpful for anticipating upcoming force changes and reducing overshoot, but it might not be optimal for accuracy and smoothness during rapid transitions. Consider adjusting the feedforward gain based on the rate of force change or using a smoother transition between setpoints.
\item Changing Setpoint Force for Trajectory Tracking\+: Feedforward Term as Next Waypoint Force\+: This strategy can aid in handling fast-\/changing waypoints, but it might lead to abrupt force variations at waypoint transitions. Consider\+: Softening the transition between waypoint forces using interpolation or a smoother feedforward update scheme. Incorporating a dynamic model-\/based feedforward term for more accurate trajectory tracking, even if it\textquotesingle{}s a simplified model.
\end{DoxyEnumerate}

Consider adjusting the feedforward gain based on the rate of force change or using a smoother transition between setpoints.

Extra Notes As force measurements are often noisy and thus their differentiation does not produce meaningful results, usually a Proportional-\/\+Integral (PI) controller is used in the force control loop and the update frequency must be high enough to avoid instabilities.

THIS IS WRONG BTW, BUT IT GIVES A GOOD IDEA OF HOW TO HANDLE IT // Compute the total force // double total\+\_\+force = control\+\_\+signal; // if (compensate\+\_\+for\+\_\+ee\+\_\+weight\+\_\+) \{ // total\+\_\+force += ee\+\_\+weight\+\_\+; // \}

Video\+: Dif PID tuning Move the board Publish motion towards the right Show the safety Show other dims Show with admittance -\/ see if they can both run Add flow diagram for controls

Rethink\+: transform\+\_\+to\+\_\+ee Rethink implementing full wrench dimensions

Todo\+: configure moveit servo and test with it Make video in sim Combine screen cast and recorded video

Note\+: When I use the force\+\_\+torque\+\_\+sensor\+\_\+broadcaster -\/ there\textquotesingle{}s a lag after starting the twist\+\_\+streaming Todo\+: add transformations and filtering and gravity compensation --add -\/ if no force is felt, keep the gain constant. Only start the controller when force is felt

December presentation Picknik controller implementaiton
\begin{DoxyItemize}
\item admittance -\/ scattered everywhere, ros2 controller, gravity comp, -\/these are the packages i wrote, so you don\textquotesingle{}t have to force control -\/ tested dif packges, implemented, drift (show this) and then implemnted my own
\item deadband, filter forces, gravity compensation (no drift) Submitted to TAHRI, humanoid coming up.
\end{DoxyItemize}

This is the contents of my next submission\+: Research -\/ spot compliance (wrench), statics turn wheel valve
\begin{DoxyItemize}
\item traj tracker, interactive screw, compare to cartesian
\end{DoxyItemize}

Has in built safety -\/ no ft, no publishing and ignores the other dimensions

different node to filter the forces

The integral term is responsible for accumulating the sum of past errors. In each iteration of the control loop, the error is added to the integral term. This accumulation helps the controller to eliminate steady-\/state errors over time. If there is a constant error over several iterations, the integral term will grow and contribute to reducing this steady-\/state error.

Tuning the Controller\+: Fine-\/tune the PID parameters to achieve the desired performance. This may involve experimenting with different proportional, integral, and derivative gains to find a balance between stability and responsiveness.

/$\ast$$\ast$$\ast$$\ast$$\ast$$\ast$$\ast$$\ast$$\ast$$\ast$$\ast$$\ast$$\ast$$\ast$$\ast$$\ast$$\ast$$\ast$$\ast$$\ast$$\ast$$\ast$$\ast$$\ast$$\ast$$\ast$$\ast$$\ast$$\ast$$\ast$$\ast$$\ast$$\ast$$\ast$$\ast$$\ast$$\ast$$\ast$$\ast$$\ast$$\ast$$\ast$$\ast$$\ast$$\ast$$\ast$$\ast$$\ast$$\ast$$\ast$$\ast$$\ast$$\ast$$\ast$$\ast$$\ast$$\ast$$\ast$$\ast$$\ast$$\ast$$\ast$$\ast$/ //$\ast$$\ast$$\ast$$\ast$$\ast$$\ast$$\ast$$\ast$$\ast$$\ast$$\ast$$\ast$$\ast$$\ast$$\ast$$\ast$$\ast$$\ast$$\ast$$\ast$$\ast$$\ast$ void Force\+Controller\+::update\+FTCB( const geometry\+\_\+msgs\+::msg\+::\+Wrench\+Stamped\+::\+Shared\+Ptr wrench)\{\} RCLCPP\+\_\+\+INFO\+\_\+\+THROTTLE(this-\/$>$get\+\_\+logger(), $\ast$this-\/$>$get\+\_\+clock(), ROS\+\_\+\+LOG\+\_\+\+THROTTLE\+\_\+\+PERIOD, \char`\"{}\+New wrench received.\char`\"{});

RCLCPP\+\_\+\+INFO\+\_\+\+THROTTLE( this-\/$>$get\+\_\+logger(), $\ast$this-\/$>$get\+\_\+clock(), ROS\+\_\+\+LOG\+\_\+\+THROTTLE\+\_\+\+PERIOD, \char`\"{}\+Current wrench\+:  force\+: \mbox{[}\%f, \%f, \%f\mbox{]}, torque\+: \mbox{[}\%f, \%f, \%f\mbox{]}\char`\"{}, cur\+\_\+wrench\+\_\+.\+wrench.\+force.\+x, cur\+\_\+wrench\+\_\+.\+wrench.\+force.\+y, cur\+\_\+wrench\+\_\+.\+wrench.\+force.\+z, cur\+\_\+wrench\+\_\+.\+wrench.\+torque.\+x, cur\+\_\+wrench\+\_\+.\+wrench.\+torque.\+y, cur\+\_\+wrench\+\_\+.\+wrench.\+torque.\+z);

/$\ast$$\ast$$\ast$$\ast$$\ast$$\ast$$\ast$$\ast$$\ast$$\ast$$\ast$$\ast$$\ast$$\ast$$\ast$$\ast$$\ast$$\ast$$\ast$$\ast$$\ast$$\ast$$\ast$$\ast$$\ast$$\ast$$\ast$$\ast$$\ast$$\ast$$\ast$$\ast$$\ast$$\ast$$\ast$$\ast$$\ast$$\ast$$\ast$$\ast$$\ast$$\ast$$\ast$$\ast$$\ast$$\ast$$\ast$$\ast$$\ast$$\ast$$\ast$$\ast$$\ast$$\ast$$\ast$$\ast$$\ast$$\ast$$\ast$$\ast$$\ast$ target\+\_\+wrench\+\_\+sub\+\_\+ = this-\/$>$create\+\_\+subscription$<$geometry\+\_\+msgs\+::msg\+::\+Wrench\+Stamped$>$( this-\/$>$get\+\_\+parameter(\char`\"{}target\+\_\+wrench\+\_\+topic\char`\"{}).as\+\_\+string(), rclcpp\+::\+Sensor\+Data\+Qo\+S(), /$\ast$\+Match the QoS policy of the publishser$\ast$/ std\+::bind(\&Force\+Controller\+::update\+Target\+Wrench\+CB, this, std\+::placeholders\+::\+\_\+1)); /////////////////////////////

///
\begin{DoxyParams}{Parameters}
{\em wrench} & void Force\+Controller\+::update\+Target\+Wrench\+CB( const geometry\+\_\+msgs\+::msg\+::\+Wrench\+Stamped\+::\+Shared\+Ptr wrench) \{ // Verify that the wrench data is valid\\
\hline
\end{DoxyParams}
if () \{ RCLCPP\+\_\+\+WARN(this-\/$>$get\+\_\+logger(), \char`\"{}\+Invalid wrench data received. Ignoring input.\char`\"{}); return; \}

target\+\_\+wrench\+\_\+ = $\ast$wrench;

//$\ast$$\ast$$\ast$$\ast$$\ast$$\ast$$\ast$$\ast$$\ast$$\ast$$\ast$$\ast$$\ast$$\ast$$\ast$$\ast$$\ast$$\ast$$\ast$$\ast$$\ast$$\ast$ RCLCPP\+\_\+\+INFO\+\_\+\+THROTTLE(this-\/$>$get\+\_\+logger(), $\ast$this-\/$>$get\+\_\+clock(), ROS\+\_\+\+LOG\+\_\+\+THROTTLE\+\_\+\+PERIOD, \char`\"{}\+New target wrench received.\char`\"{});

RCLCPP\+\_\+\+INFO\+\_\+\+THROTTLE( this-\/$>$get\+\_\+logger(), $\ast$this-\/$>$get\+\_\+clock(), ROS\+\_\+\+LOG\+\_\+\+THROTTLE\+\_\+\+PERIOD, \char`\"{}\+Target wrench\+:  force\+: \mbox{[}\%f, \%f, \%f\mbox{]}, torque\+: \mbox{[}\%f, \%f, \%f\mbox{]}\char`\"{}, target\+\_\+wrench\+\_\+.\+wrench.\+force.\+x, target\+\_\+wrench\+\_\+.\+wrench.\+force.\+y, target\+\_\+wrench\+\_\+.\+wrench.\+force.\+z, target\+\_\+wrench\+\_\+.\+wrench.\+torque.\+x, target\+\_\+wrench\+\_\+.\+wrench.\+torque.\+y, target\+\_\+wrench\+\_\+.\+wrench.\+torque.\+z); \} 